% Fuentes: EXP-010 (EXP-010_bs256_sqrt_60ep), EXP-006 (capacidad), DEC-008.
\section{Punto de partida}
El punto de partida es el modelo final de la entrega anterior (Cuadro~\ref{tab:baseline}): \res{phase0-hidden}
[\texttt{Dense} $\rightarrow$ LeakyReLU] sobre imágenes de \res{phase0-res}~$\times$~\res{phase0-res}
(\res{phase0-params} parámetros), entrenado con Adam y \emph{batch} de \res{phase0-batch}, con \emph{EarlyStopping} y
\emph{ReduceLROnPlateau} sobre el \emph{mean recall} de validación. Tres observaciones de la entrega anterior apuntan
a sobreajuste:
\begin{itemize}
  \item \textbf{Brecha.} Su \emph{accuracy} de entrenamiento (\res{phase0-train-acc}) supera a la de validación
  (\res{phase0-val-acc}) en \res{phase0-gap-acc}~pp.
  \item \textbf{Capacidad.} Aumentar el número de parámetros solo aumentaba esa brecha: una red de tres capas
  (1024--512--256, \res{big3layer-params} parámetros) llegaba a una \emph{accuracy} de entrenamiento de
  \res{big3layer-train-acc} con el mismo \emph{mean recall} de validación (\res{big3layer-mean-recall}).
  \item \textbf{Parada.} Tras la mejor época, la \emph{accuracy} de entrenamiento seguía subiendo y la de validación
  no: \emph{EarlyStopping} paraba porque la red solo memorizaba.
\end{itemize}
La regularización restringe el modelo para que generalice mejor, aunque empeore en entrenamiento~\cite{goodfellow}.
Si la brecha es varianza (sobreajuste), regularizar debería cerrarla \emph{subiendo} la validación; si lo que limita es
el sesgo del modelo, la cerrará bajando el entrenamiento, sin ganancia en validación. Los experimentos de esta entrega
distinguen los dos casos.

\begin{table}[h]
  \centering
  \safetable{generated/tables/baseline_summary}
  \caption{Modelo de partida (final de la entrega 1) y sus métricas (media $\pm$ desviación típica de \res{phase0-seeds}
  semillas).}
  \label{tab:baseline}
\end{table}
